\section{Circuits, gate sets, and the Hilbert--Schmidt distance}
\label{sec:circuits}

This section fixes the vocabulary shared by the whole Clifford+T development:
what a circuit is, which gate sets appear, and which distance measures
approximation.  Everything here is formalized and proved.

\subsection{Circuits and synthesis}

\begin{definition}
  \label{def:circuit_matrix}
  \lean{TwoControl.Clifford.Universal.circuitMatrix}
  \leanok
  A circuit on $N$-dimensional state space is a list of $N \times N$ matrices,
  and its matrix is the product taken left to right:
  \[
    \operatorname{circuitMatrix}[\,] = I, \qquad
    \operatorname{circuitMatrix}(g :: gs)
      = g \cdot \operatorname{circuitMatrix}(gs).
  \]
\end{definition}

\begin{definition}
  \label{def:circuit_over}
  \lean{TwoControl.Clifford.Universal.CircuitOver}
  \leanok
  For a predicate $\mathcal{G}$ on $N \times N$ matrices,
  $\operatorname{CircuitOver} \mathcal{G}\, gs$ says every factor of $gs$
  satisfies $\mathcal{G}$.  Concatenation of circuits multiplies their
  matrices, and a concatenation lies over $\mathcal{G}$ iff both parts do.
\end{definition}

\begin{definition}
  \label{def:global_phase}
  \lean{TwoControl.Clifford.GlobalPhaseEquivalent}
  \leanok
  Two matrices are \emph{globally phase equivalent}, written $U \sim V$, when
  $U = z \cdot V$ for some scalar $z$ with $\lvert z \rvert = 1$.  This
  relation is reflexive and transitive, is compatible with multiplication, and
  can be transported through the local tensor factors $A \mapsto A \otimes I$
  and $A \mapsto I \otimes A$.
\end{definition}

\begin{definition}
  \label{def:synthesizes}
  \lean{TwoControl.Clifford.Universal.Synthesizes}
  \leanok
  $\operatorname{Synthesizes} \mathcal{G}\, U$ asserts that some circuit over
  $\mathcal{G}$ has matrix exactly $U$.
\end{definition}

\begin{definition}
  \label{def:synthesizes_phase}
  \lean{TwoControl.Clifford.Universal.SynthesizesUpToGlobalPhase}
  \leanok
  $\operatorname{SynthesizesUpToGlobalPhase} \mathcal{G}\, U$ asserts that
  some circuit over $\mathcal{G}$ has matrix globally phase equivalent to $U$.
  This is the shape in which the paper's exact synthesis results hold: the
  one-qubit base case produces a residual scalar that the allowed gate set
  cannot itself realize.
\end{definition}

\subsection{Embedding one- and two-qubit gates into a register}

\begin{definition}
  \label{def:embedded_one}
  \lean{TwoControl.Clifford.Universal.IsEmbeddedOneQubitGate}
  \leanok
  A \emph{placement} records a choice of wire inside an $n$-qubit register
  together with the reindexing that identifies
  $\mathbb{C}^{2^n}$ with $\mathbb{C}^{2^a} \otimes \mathbb{C}^{2}
  \otimes \mathbb{C}^{2^b}$.  Given a one-qubit gate $g$, the predicate
  $\operatorname{IsEmbeddedOneQubitGate} n\, g\, E$ says $E$ is $g$ acting on
  one wire of an $n$-qubit register (possibly as one half of a two-wire
  placement).
\end{definition}

\begin{definition}
  \label{def:embedded_two}
  \lean{TwoControl.Clifford.Universal.IsEmbeddedTwoQubitGate}
  \leanok
  Likewise $\operatorname{IsEmbeddedTwoQubitGate} n\, V\, E$ says $E$ is the
  two-qubit gate $V$ acting on an ordered pair of distinct wires of an
  $n$-qubit register.  Embedded gates of either kind are unitary whenever the
  gate is.
\end{definition}

\subsection{The gate sets}

\begin{definition}
  \label{def:gate_cliffordrz}
  \lean{TwoControl.Clifford.Universal.CliffordRzGate}
  \leanok
  The \emph{paper gate set} of \texttt{universal\_new\_gates.tex} is
  \[
    \mathcal{G}_{\mathrm{Clifford}+R_z}
      = \{\, C(X),\; H,\; S,\; S^\dagger,\; R_z(\theta) \,\},
  \]
  each embedded on wires of an $n$-qubit register.  This is the target of
  exact synthesis (Theorem~\ref{thm:clifford_rz_universal}).
\end{definition}

\begin{definition}
  \label{def:gate_cliffordtrz}
  \lean{TwoControl.Clifford.Universal.CliffordTRzGate}
  \leanok
  The \emph{intermediate gate set} is
  $\{\, C(X),\, H,\, T,\, R_z(\theta) \,\}$: the Clifford phases $S, S^\dagger$
  have been rewritten as powers of $T$, but $R_z(\theta)$ is still exact.
\end{definition}

\begin{definition}
  \label{def:gate_cliffordt}
  \lean{TwoControl.Clifford.Universal.CliffordTGate}
  \leanok
  The \emph{final gate set} is $\{\, C(X),\, H,\, T \,\}$, obtained by
  dropping $R_z(\theta)$.  This is the gate set of the main theorem.
\end{definition}

\begin{definition}
  \label{def:gate_easy}
  \lean{TwoControl.Clifford.Universal.EasyGate}
  \leanok
  The \emph{easy gate set} allows an arbitrary embedded two-qubit unitary
  together with embedded $H$, $S$, $S^\dagger$, and $R_z(\theta)$.  It is a
  weakening of Definition~\ref{def:gate_cliffordrz}: every paper gate is an
  easy gate, since $C(X)$ is in particular a unitary two-qubit gate.  It
  belongs to the older decomposition route of
  Section~\ref{sec:easy_gate_layer} and is not used by the main theorem.
\end{definition}

\begin{definition}
  \label{def:ht_primitive}
  \lean{TwoControl.Clifford.Universal.OneQubitHTPrimitive}
  \leanok
  A two-constructor type naming the one-qubit primitives $H$ and $T$, with an
  evaluation map to $2 \times 2$ matrices and a circuit-matrix function
  $\operatorname{oneQubitHTCircuitMatrix}$.  A parallel type
  $\operatorname{TwoQubitCliffordTPrimitive}$ packages local two-qubit
  Clifford+T words.  Both come with lifting maps that turn a local circuit into
  an $n$-qubit circuit over $\operatorname{CliffordTGate}$, together with the
  lemma that lifting commutes with taking the circuit matrix.  These are what
  let the one-qubit output of Lemma 12 be spliced into an $n$-qubit circuit.
\end{definition}

\subsection{The Hilbert--Schmidt distance}

\begin{definition}
  \label{def:hs_distance}
  \lean{TwoControl.Clifford.Universal.hsDistance}
  \leanok
  For $N \times N$ matrices $U, V$,
  \[
    d(U,V) \;=\; \sqrt{\,1 - \frac{\lvert \operatorname{Tr}(U^\dagger V)
      \rvert^2}{N^2}\,}.
  \]
  This is the distance in which the paper states approximation, and in which
  the main theorem is proved.
\end{definition}

The lemmas that follow are the paper's distance toolkit.  All are proved.

\begin{lemma}
  \label{lem:hs_self}
  \lean{TwoControl.Clifford.Universal.hsDistance_self}
  \leanok
  For unitary $U$, $d(U,U) = 0$.
\end{lemma}

\begin{proof}
  \leanok
  \uses{def:hs_distance}
  $\operatorname{Tr}(U^\dagger U) = \operatorname{Tr}(I_N) = N$, so the
  radicand is $1 - N^2/N^2 = 0$.
\end{proof}

\begin{lemma}
  \label{lem:hs_phase}
  \lean{TwoControl.Clifford.Universal.hsDistance_eq_zero_of_globalPhaseEquivalent}
  \leanok
  If $U$ and $V$ are unitary and $U \sim V$, then $d(U,V) = 0$.
\end{lemma}

\begin{proof}
  \leanok
  \uses{def:hs_distance, def:global_phase, lem:hs_self}
  Write $U = z V$ with $\lvert z \rvert = 1$.  Then
  $\operatorname{Tr}(U^\dagger V) = \bar{z}\operatorname{Tr}(V^\dagger V)
   = \bar{z} N$, whose modulus is $N$.
\end{proof}

\begin{lemma}
  \label{lem:hs_trace_ineq}
  \lean{TwoControl.Clifford.Universal.trace_inequality}
  \leanok
  For $N \times N$ unitaries $U, V$,
  \[
    \sqrt{1 - \frac{\lvert\operatorname{Tr}(UV)\rvert^2}{N^2}}
      \;\le\;
    \sqrt{1 - \frac{\lvert\operatorname{Tr}(U)\rvert^2}{N^2}}
      +
    \sqrt{1 - \frac{\lvert\operatorname{Tr}(V)\rvert^2}{N^2}}.
  \]
\end{lemma}

\begin{proof}
  \leanok
  This is the inequality the paper cites from Wang and Zhang (1994).  It is
  proved here rather than assumed: normalising the three matrices to unit
  vectors in the Hilbert--Schmidt inner product turns the claim into the
  statement that $x \mapsto \sqrt{1 - \lvert\langle x, y\rangle\rvert^2}$
  satisfies a triangle inequality on the unit sphere, which is discharged by
  an inner-product-space computation.
\end{proof}

\begin{lemma}
  \label{lem:hs_mul}
  \lean{TwoControl.Clifford.Universal.hsDistance_mul_le}
  \leanok
  For unitaries $U_1, U_2, V_1, V_2$,
  \[
    d(U_1 U_2, V_1 V_2) \le d(U_1, V_1) + d(U_2, V_2).
  \]
\end{lemma}

\begin{proof}
  \leanok
  \uses{lem:hs_trace_ineq, lem:hs_rearrange}
  Rewrite $\operatorname{Tr}\!\big((U_1U_2)^\dagger V_1 V_2\big)$ cyclically as
  $\operatorname{Tr}\!\big((V_1^\dagger U_1)(V_2 U_2^\dagger)\big)$ using
  Lemma~\ref{lem:hs_rearrange}, then apply
  Lemma~\ref{lem:hs_trace_ineq} to the two factors.
\end{proof}

\begin{lemma}
  \label{lem:hs_circuit_sum}
  \lean{TwoControl.Clifford.Universal.hsDistance_circuitMatrix_le_sum}
  \leanok
  If two circuits have equal length and unitary factors $U_i$ and $V_i$, then
  \[
    d\big(\operatorname{circuitMatrix}(U_\bullet),
          \operatorname{circuitMatrix}(V_\bullet)\big)
      \;\le\; \sum_i d(U_i, V_i).
  \]
\end{lemma}

\begin{proof}
  \leanok
  \uses{lem:hs_mul, def:circuit_matrix}
  Induct on the length, applying Lemma~\ref{lem:hs_mul} at each
  multiplication.
\end{proof}

\begin{lemma}
  \label{lem:hs_rearrange}
  \lean{TwoControl.Clifford.Universal.hsDistance_product_rearrange}
  \leanok
  $d(U_1 U_2, V_1 V_2) = d(V_1^\dagger U_1, V_2 U_2^\dagger)$.
\end{lemma}

\begin{proof}
  \leanok
  \uses{def:hs_distance}
  Both sides have the same trace argument after cyclic rearrangement, and the
  definition of $d$ depends only on that trace.
\end{proof}

\begin{lemma}
  \label{lem:hs_identity}
  \lean{TwoControl.Clifford.Universal.hsDistance_one_mul_mul}
  \leanok
  $d(I, V_1 V_2) = d(V_1^\dagger, V_2)$, and symmetrically
  $d(U_1 U_2, I) = d(U_1, U_2^\dagger)$.
\end{lemma}

\begin{proof}
  \leanok
  \uses{lem:hs_rearrange}
  Specialize Lemma~\ref{lem:hs_rearrange}.
\end{proof}

\begin{lemma}
  \label{lem:hs_dagger}
  \lean{TwoControl.Clifford.Universal.hsDistance_conjTranspose_symm}
  \leanok
  $d(U,V) = d(V^\dagger, U^\dagger)$.
\end{lemma}

\begin{proof}
  \leanok
  \uses{def:hs_distance}
  $\operatorname{Tr}\!\big((V^\dagger)^\dagger U^\dagger\big)
    = \operatorname{Tr}\!\big((U^\dagger V)^\dagger\big)$, and conjugation does
  not change the modulus of a trace.
\end{proof}

\begin{lemma}
  \label{lem:hs_triangle_one}
  \lean{TwoControl.Clifford.Universal.hsDistance_le_hsDistance_one_add}
  \leanok
  For unitaries $U, V$: $d(U,V) \le d(U,I) + d(I,V)$.
\end{lemma}

\begin{proof}
  \leanok
  \uses{lem:hs_trace_ineq, lem:hs_identity}
  This is the paper's shape of the trace inequality; it follows from
  Lemma~\ref{lem:hs_trace_ineq} after rewriting both distances against the
  identity using Lemma~\ref{lem:hs_identity}.
\end{proof}
